\scanpage{41}
% archive sheet number: 40
b) \uline{Filtration}s\deleted{\unsure{ass}} \uline{associées aux cristaux de Dieudonné}. Revenons au cas d'un $S$ général (sur lequel $p$ soit loc. nilp.). On défini\deleted{\unsure{t}}ra alors, pour tout $G\handwritten{\in}\mathrm{Ob}\,\mathrm{BT}(S)$, une filtration \deleted{en deux crans} d\deleted{e}\unsure{u|e} Mo\deleted{\unsure{d|D}}\unsure{D|d}ule localement libre $\underline{D}^{\handwritten{*}}(G)_S$ sur $S$ par un sous-Module localement facteur direct (donc localement libre de t.f.) $\inserted{\handwritten{\underline{\mathrm{Fil}}^1=}}\underline{\handwritten{\omega}}_G$, donnant lieu à une suite exacte
\[
0\handwritten{\longrightarrow}\underline{\handwritten{\omega}}_G\handwritten{\longrightarrow}\underline{D}^{\handwritten{*}}(G)\handwritten{\longrightarrow}\underline{t}_{G^{\handwritten{*}}}\handwritten{\longrightarrow}0\quad,
\tag*{\handwritten{($\unsure{*|\delta}$)}}
\]
où en fait $\underline{\handwritten{\omega}}_G$ est le module des différentielles invariantes sur le groupe formel $\bar{G}\handwritten{\subset}G$ associé à $G$ (qui sera construit plus bas), et $\underline{t}_G=\underline{\handwritten{\omega}}_{G^{\handwritten{*}}}^{\handwritten{\vee}}$ est l'Algèbre de Lie associée au gro\deleted{\unsure{u|v}}upe de BT dual $G^{\handwritten{*}}$ de $G$ (définition à reporter aux généralités sur les goupes de BT). Cette \deleted{filtr} suite exacte est fonctorielle \inserted{en $G$} \deleted{et} compatible aux images inverses par un $S'\handwritten{\longrightarrow}S$.

\margin{\handwritten{(*) Sans doute inutile si $S_{\unsure{o|0}}$ \unsure{parfait|parfaite}.}}
